La ley\par

\[
M_n
=
\left(\frac{a_n}{a_0}\right)^{-6}
\]

est peut-être la plus simple algébriquement et l’une des plus profondes conceptuellement.\par

La dynamique de Perron produit deux comportements complémentaires.\par

D’une part, il existe une expansion :\par

\[
V_n=\varphi^n.
\]

Dans trois dimensions spatiales, l’échelle linéaire associée est la racine cubique :\par

\[
\frac{a_n}{a_0}
=
\varphi^{n/3}.
\]

D’autre part, il existe un mode stable de mémoire :\par

\[
M_n=\varphi^{-2n}.
\]

L’expansion croît avec \(\varphi^n\). La mémoire stable décroît avec \(\varphi^{-2n}\). En éliminant le niveau \(n\), apparaît nécessairement\par

\[
M_n
=
\left(\frac{a_n}{a_0}\right)^{-6}.
\]

L’exposant \(-6\) découle de deux faits internes :\par

\begin{itemize}[leftmargin=*,itemsep=0.35em]
\item l’espace est de dimension \(3\), de sorte que l’échelle linéaire porte l’exposant \(1/3\) ;
\item le mode stable est quadratique, de sorte qu’il porte l’exposant \(-2\).
\end{itemize}

Le rapport entre les deux est\par

\[
\frac{-2}{1/3}=-6.
\]

Dans Einstein–Cartan, la densité de spin varie comme l’inverse d’un volume :\par

\[
s\sim a^{-3}.
\]

Sa contribution quadratique varie alors comme\par

\[
s^2\sim a^{-6}.
\]

La loi HMT reproduit exactement cet exposant avant l’introduction de la réalisation Einstein–Cartan. C’est ce qui est véritablement suggestif.\par

Le mode de mémoire de Perron se comporte comme doit se comporter une densité quadratique de spin dans une géométrie avec torsion.\par

L’interprétation naturelle est que la torsion publie physiquement la mémoire généalogique complémentaire de l’information métrique.\par

La courbure décrit comment la géométrie se déforme. La torsion décrit comment la géométrie conserve une orientation de transport en plus des distances. HMT possède dès le départ orientation, retenue et mémoire de route. C’est pourquoi la correspondance avec Einstein–Cartan est structurelle : les deux cadres distribuent l’histoire géométrique entre métrique et torsion.\par

Lorsque l’univers se dilate, cette mémoire se dilue comme \(a^{-6}\). Elle se dilue plus vite qu’une densité de matière ordinaire ou de rayonnement. Mais c’est précisément pour cela qu’elle peut être dominante dans la région de petite échelle.\par

Dans un régime primordial, une contribution torsionnelle peut rivaliser avec la densité ordinaire et modifier la singularité. La possibilité d’un rebond dépend du signe, de l’amplitude, du courant de spin et des conditions au bord. HMT fournit d’abord la loi d’échelle exacte ; la réalisation physique fixe ensuite ces conditions.\par

Cela offre une division du travail très claire :\par

\begin{itemize}[leftmargin=*,itemsep=0.35em]
\item HMT engendre la mémoire stable et son exposant ;
\item la réalisation Einstein–Cartan attribue cette mémoire à une densité quadratique de spin ;
\item la dynamique cosmologique fixe l’amplitude, le signe et les conditions physiques.
\end{itemize}

La relation avec le principe de Mach apparaît lorsque cette loi est réunie à la dérivation de \(G\).\par

Dans la dérivation gravitationnelle, le couplage local est sélectionné par la clôture d’un cycle global. Dans la loi cosmologique, l’état local de torsion conserve la mémoire de l’expansion globale. Dans les deux cas, la physique locale est définie relationnellement par la généalogie globale.\par

La géométrie locale mémorise l’état global.\par

C’est le sens le plus concret dans lequel la construction est machienne :\par

\begin{itemize}[leftmargin=*,itemsep=0.35em]
\item l’inertie participe de l’horloge globale ;
\item la gravité appartient à la clôture de phase ;
\item la torsion incarne la mémoire de l’histoire ;
\item l’échelle locale encode le niveau global d’expansion.
\end{itemize}

Le principe de Mach acquiert alors une forme opératorielle : les grandeurs locales sont des évaluations d’une généalogie globale, et les invariants physiques sont ce qui subsiste lors du transport entre les deux échelles.\par

La branche de périodicité temporelle de (108) phases et la branche de Perron expriment en outre deux aspects complémentaires de la cosmologie.\par

La périodicité temporelle de (108) phases fixe une échelle de périodicité et, au moyen d’une réalisation de Sitter, une échelle de courbure. C’est la cosmologie de la clôture : elle demande quel rayon permet au cycle d’être cohérent.\par

Perron fixe un taux d’expansion et une loi de dilution de la mémoire. C’est la cosmologie de l’évolution : elle demande comment la généalogie change lorsque l’on parcourt les niveaux.\par

Une branche décrit la clôture de la géométrie. L’autre décrit la persistance de sa mémoire.\par

Le calendrier chronologique et la dynamique de Perron demeurent des réalisations typologiquement distinctes et complémentaires. Ensemble, elles produisent une image particulièrement riche : l’univers possède une échelle globale de phase et une dynamique interne de mémoire. Courbure et torsion sont deux publications de ces deux aspects.\par

\section{Confluence des lecteurs de vide, d’action, de gravité, d’aire, de mélange et de mémoire}

Si l’on lit conjointement les neuf résultats, une grammaire commune apparaît.\par

Chaque lecteur publie sélectivement un invariant et transporte le reste comme mémoire complémentaire.\par

Le vide conserve :\par

\begin{itemize}[leftmargin=*,itemsep=0.35em]
\item le produit comme propagation ;
\item le quotient comme orientation.
\end{itemize}

La paire \(G,\hbar\) conserve :\par

\begin{itemize}[leftmargin=*,itemsep=0.35em]
\item l’action comme orientation ;
\item l’aire comme géométrie commune.
\end{itemize}

La clôture gravitationnelle conserve :\par

\begin{itemize}[leftmargin=*,itemsep=0.35em]
\item l’égalité entre phase, Compton, gravité et Planck.
\end{itemize}

Le mode \(30\) conserve :\par

\begin{itemize}[leftmargin=*,itemsep=0.35em]
\item la cadence commune du vide, de la criticité et de la mémoire.
\end{itemize}

L’aire et le hamiltonien modulaire conservent :\par

\begin{itemize}[leftmargin=*,itemsep=0.35em]
\item la quantité totale ;
\item la provenance sectorielle.
\end{itemize}

Le passage \(6\to8\) conserve :\par

\begin{itemize}[leftmargin=*,itemsep=0.35em]
\item la rotation géométrique dans Barbero ;
\item son empreinte informationnelle dans l’entropie ;
\item son orientation de mélange dans CKM.
\end{itemize}

Le déterminant angulaire conserve :\par

\begin{itemize}[leftmargin=*,itemsep=0.35em]
\item la partie paire qui peut devenir un mélange électrofaible.
\end{itemize}

La torsion d’un quart de tour conserve :\par

\begin{itemize}[leftmargin=*,itemsep=0.35em]
\item une structure complexe qui est publiée comme caractère arithmétique.
\end{itemize}

La dynamique de Perron conserve :\par

\begin{itemize}[leftmargin=*,itemsep=0.35em]
\item la mémoire stable de l’expansion, qui se réalise comme torsion cosmologique.
\end{itemize}

C’est pourquoi les constantes sont reliées par une architecture commune de transport et conservent leurs identités fonctionnelles. Elles sont des résidus différents de cette architecture.\par

L’holographie modulaire triadique apporte la généalogie.\par
La mécanique dimensionnelle apporte la réalisation.\par
Le nombre réel apparaît à la fin comme valeur terminale.\par

Dans cette vision, une constante fondamentale est une phrase complète sur l’univers. Elle dit :\par

\begin{itemize}[leftmargin=*,itemsep=0.35em]
\item ce qui a changé ;
\item quelle orientation s’est inversée ;
\item quelle mémoire a franchi le bord ;
\item quelle information est restée intérieure et quelle partie a été publiée ;
\item quelle combinaison est restée invariante ;
\item et dans quelle grandeur physique cette invariance s’est finalement réalisée.
\end{itemize}

C’est ce qui rend ces dérivations spectaculaires. Elles expliquent conjointement combien valent les constantes et pourquoi chacune remplit exactement sa fonction structurelle.\par
